\documentclass[11pt,a4paper]{article}

\newcommand{\ReportShortTitle}{From informal claims to checked theorems}
\newcommand{\ReportHeaderRight}{Process history}
\usepackage{reports/ipnpcns-report}

\hypersetup{
  pdftitle={From Informal Claims to Checked Theorems: An Evidence-Based History of the IPNPCNS Formalization},
  pdfauthor={IPNPCNS formalization project},
  pdfsubject={Evidence-based history of a Lean formalization and verification project},
  pdfkeywords={Lean, formalization, proof development, project history, convex cones, verification evidence}
}

\newcommand{\artifact}[1]{\texttt{\detokenize{#1}}}
\newcommand{\commit}[1]{\texttt{#1}}
\newcommand{\ev}[1]{%
  \hyperref[ev:#1]{\textcolor{ReportBlue}{\textsf{\scriptsize[#1]}}}}
\newcommand{\evidenceitem}[2]{%
  \item[\textbf{#1}]\phantomsection\label{ev:#1}#2}

\newtcolorbox{turningpoint}[1][]{
  enhanced,
  breakable,
  colback=ReportWarm,
  colframe=ReportGold,
  boxrule=0.8pt,
  arc=1.2mm,
  left=2.5mm,right=2.5mm,top=2mm,bottom=2mm,
  title={#1},
  fonttitle=\bfseries,
  coltitle=ReportInk
}

\newtcolorbox{evidencebox}[1][]{
  enhanced,
  breakable,
  colback=white,
  colframe=ReportCyan,
  boxrule=0.65pt,
  arc=1.2mm,
  left=2.5mm,right=2.5mm,top=2mm,bottom=2mm,
  title={#1},
  fonttitle=\bfseries,
  coltitle=ReportNavy
}

\begin{document}

\hypersetup{pageanchor=false}
\begin{titlepage}
  \thispagestyle{empty}
  \begin{tikzpicture}[remember picture,overlay]
    \fill[ReportNavy] (current page.north west)
      rectangle ([yshift=-93mm]current page.north east);
    \fill[ReportGold] ([yshift=-93mm]current page.north west)
      rectangle ([yshift=-97mm]current page.north east);
    \fill[ReportWarm] ([yshift=27mm]current page.south west)
      rectangle (current page.south east);
  \end{tikzpicture}

  \vspace*{13mm}
  {\color{white}\sffamily\bfseries\fontsize{29}{34}\selectfont
    From Informal Claims\par
    to Checked Theorems\par}
  \vspace{7mm}
  {\color{white}\sffamily\Large
    An evidence-based history of the IPNPCNS formalization\par}

  \vfill

  \begin{tcolorbox}[
    colback=white,
    colframe=ReportGold,
    boxrule=0.8pt,
    arc=1.5mm,
    left=4mm,right=4mm,top=3mm,bottom=3mm
  ]
  \begin{tabularx}{\linewidth}{@{}L{39mm}Y@{}}
    \textbf{Subject} &
      How a 60-page interdisciplinary paper became a layered Lean development,
      and how proof failures, specification boundaries, and audits shaped the result \\
    \textbf{Primary record} &
      Wotan task logs, project-control decisions, Git states, compiler and build
      output, source audit, and a checksum-bounded Codex session trace \\
    \textbf{Final proof corpus} &
      31 Lean modules, 8,367 lines, 439 theorem or lemma declarations, and no
      admitted proofs or project-specific axioms \\
    \textbf{Historical stance} &
      Causal reconstruction with explicit evidence classes; not a task diary and
      not a claim that the final architecture was inevitable \\
    \textbf{Date} & August 2026
  \end{tabularx}
  \end{tcolorbox}

  \vspace{12mm}
  {\large\color{ReportNavy}\textbf{IPNPCNS formalization project}}\par
  \vspace{2mm}
  {\color{ReportGray}
    A companion to the separate scientific report on the verified mathematics.}
\end{titlepage}

\hypersetup{pageanchor=true}
\pagenumbering{roman}

\begin{abstract}
This report reconstructs how the mathematical claims in a paper on neural
population computation were converted into a machine-checked Lean development.
Its central thesis is that the work did not proceed by translating formulas in
order. It proceeded by repeated \emph{contract refinement}: an informal claim was
given an explicit domain, closure convention, quantifier structure, and verification
boundary; Lean and mathlib then supplied counterpressure through type errors,
missing infrastructure, counterexamples, and proof obligations; the resulting
contract was either proved, split, specialized, or retained as a visible premise.

The history is unusually well documented. Git preserves 24 curated checkpoints
through the first report. Wotan logs and project-control records preserve declared
scope and steering decisions. A 50 MB, checksum-bounded Codex session prefix also
preserves the dense work between commits: full additions and diffs for discarded
scratch proofs, compiler diagnostics, theorem-name probes, and migrations into the
permanent modules. This record permits a distinction between ordinary elaboration
failures, genuine mathematical obstructions, architectural changes of proof, and
late verification or typesetting corrections.

Several episodes were decisive. A coordinate-free Moreau/double-rejection pilot
validated the general Hilbert-space architecture. Rank deficiency disproved strict
coefficient contraction and forced a Fejer/compactness convergence proof. Missing
polyhedral-face, Moore--Penrose, and Hilbert--Schmidt APIs led respectively to
sufficient active-set certificates, an effective-map pseudoinverse, and a compact
Volterra construction by finite-rank approximation. Analytic and probabilistic
claims were first isolated behind named interfaces; some were later discharged by
Lipschitz, exponential-error, or independent-product theorems, while biological
learning premises remained explicit. A final audit against the LaTeX source found no
material mathematical discrepancy in the claimed theorem set.

The resulting lesson is methodological rather than celebratory. Fast autonomous
formalization remains credible only when theorem contracts, discarded attempts,
assumption boundaries, build evidence, and source conformance are all preserved as
different forms of evidence.
\end{abstract}

\begingroup
\small
\setlength{\parskip}{0pt}
\tableofcontents
\endgroup
\clearpage

\pagenumbering{arabic}

\section{The history this report tells}
\label{sec:purpose}

There are two legitimate accounts of a formalization, and confusing them weakens
both. The companion scientific report asks what the model says, which theorems were
proved, which assumptions they require, and what machine checking certifies. This
report asks a different question: how did the final theorem set, architecture, and
standard of evidence come into being? The distinction was made an explicit closure
requirement before the proof phase ended. \ev{E23}\ev{E24}

The answer cannot be a list of tasks. A chronological list would show that one file
followed another, but not why finite subspaces preceded pseudoinverses, why a strict
NNLS contraction bifurcated into two convergence theorems, or why a generic
convolution claim became a concrete causal Volterra theorem. The useful unit of
history is therefore a \emph{causal development arc}: a claim, the pressure applied
to its first formal contract, the decision that followed, and the mathematical
object that survived.

\begin{keyresult}[Narrative thesis]
The formalization advanced through progressive theorem-contract refinement. Its
most important results came not only from completing planned proofs but from learning
which distinctions the paper's prose left implicit: pointwise image versus closed
conic hull, deterministic implication versus probability theorem, prediction
uniqueness versus coefficient uniqueness, a concrete boundary convention versus a
generic convolution, and mathematical realization versus biological validity.
\end{keyresult}

This pattern resembles the conjecture--criticism--revision cycle familiar from
mathematical practice \cite{lakatos1976}, but here the criticism was partly mechanized. Lean did not
decide research scope. It did make type, domain, dependency, and logical gaps
impossible to ignore. Project control supplied the complementary decision layer:
whether a gap should be repaired, exposed as a premise, decomposed, or left outside
the theorem set. \ev{E02}\ev{E07}\ev{E10}\ev{E14}

\begin{figure}[htbp]
\centering
\begin{tikzpicture}[
  node distance=8mm and 9mm,
  every node/.style={font=\small,align=center},
  box/.style={draw=ReportBlue,rounded corners=1mm,fill=ReportPale,
    minimum width=28mm,minimum height=12mm,inner sep=2mm},
  boundary/.style={draw=ReportGold,rounded corners=1mm,fill=ReportWarm,
    minimum width=28mm,minimum height=12mm,inner sep=2mm},
  arrow/.style={-{Latex[length=2.2mm]},thick,draw=ReportGray}
]
  \node[box] (claim) {informal claim\\and context};
  \node[box,right=of claim] (contract) {explicit theorem\\contract};
  \node[box,right=of contract] (pressure) {Lean and library\\counterpressure};
  \node[boundary,below=of pressure] (decision) {prove, split,\\specialize, or expose};
  \node[box,left=of decision] (kernel) {kernel-checked\\endpoint};
  \node[box,left=of kernel] (audit) {signature, axiom,\\and source audit};
  \draw[arrow] (claim) -- (contract);
  \draw[arrow] (contract) -- (pressure);
  \draw[arrow] (pressure) -- (decision);
  \draw[arrow] (decision) -- (kernel);
  \draw[arrow] (kernel) -- (audit);
  \draw[arrow] (audit.west) -- ++(-6mm,0) |- node[pos=0.25,left,
    text=ReportGray,font=\scriptsize]{reopen if needed} (claim.south);
\end{tikzpicture}
\caption{The recurring development cycle. The arrows describe logical
counterpressure and revision, not a rigid procedural sequence.}
\label{fig:cycle}
\end{figure}

\subsection{A result without inevitability}

The final tree is clean: named modules, stable public endpoints, a complete build,
and a source audit. That cleanliness is an endpoint property. It should not be read
backward as evidence that the path was direct. The session trace records hundreds of
intermediate patches, temporary theorem probes, failed rewrites, typeclass problems,
and deleted scratch files. Several permanent modules were literally moved from
multi-stage scratch developments only after their statements and proof routes had
stabilized. \ev{E26}

Nor did every failure have equal significance. An unavailable theorem name is an API
lookup problem. A deterministic heartbeat timeout may indicate an elaboration shape
that should be redesigned. A nonzero kernel vector fixed by an iteration is a
mathematical counterexample to strict contraction. The history becomes informative
only when these are separated.

\section{Evidence and reconstruction}
\label{sec:evidence}

\subsection{Five kinds of record}

The reconstruction uses five provenance classes. No single one is treated as a
complete account.

\begin{table}[htbp]
\centering
\caption{Historical evidence classes and their proper use.}
\label{tab:evidence-classes}
\begin{tabularx}{\linewidth}{@{}C{11mm}L{35mm}Y@{}}
\toprule
\textbf{Class} & \textbf{Record} & \textbf{What it supports} \\
\midrule
C & Contemporaneous task record &
Declared context, acceptance criteria, implementation summary, obstacle, and
verification evidence at task closeout. \\
D & Steering decision &
What evidence and uncertainty were considered at a gate, and why the next direction
was selected. \\
G & Git state &
Durable file content, order of curated checkpoints, commit timestamps, and the diff
between accepted states. \\
R & Retrospective audit &
Later source conformance, coverage classification, trust review, and report QA. It
corroborates the final state but must not be projected backward as prior knowledge. \\
S & Session trace &
Tool-mediated patches, scratch evolution, compiler feedback, and contemporaneous
updates between commits. It does not reveal unexpressed cognition. \\
\bottomrule
\end{tabularx}
\end{table}

The evidence identifiers in the margin-like blue brackets refer to the ledger in
\cref{app:evidence-ledger}. That ledger, in turn, points to the repository artifacts
and the separately bounded session audit. Material historical propositions in this
report are tied to at least one such identifier.

\subsection{The session trace changes the record}

An initial evidence review contained an overly cautious statement: no record was
known to preserve every abandoned tactic or scratch proof. The owner pointed to the
full Codex JSONL session under \artifact{~/.codex/}. Inspection changed the
historiographical conclusion. \ev{E26}

The audited prefix contains 427 successful patch-completion events. Of these, 174
changes concern 33 scratch, probe, or axiom-audit files: 34 additions, 128 updates,
and 12 deletions. Added files carry full content; updates carry unified diffs; command
outputs retain the Lean diagnostics. The 33 patches to \artifact{ScratchT17.lean}
and 26 patches to \artifact{ScratchT18.lean}, for example, survive even though those
paths do not exist in the final tree. The trace also records their moves into the
wavelet and Volterra modules. \ev{E26}

\begin{evidencebox}[Session-trace boundary]
The main session prefix is fixed by a SHA-256 digest and a 6,276-record cutoff. The
audit reads observable messages, tool calls, patches, and outputs. It deliberately
does not reproduce model-internal reasoning records or copy the raw 50 MB log into
the repository. Thus the report can reconstruct expressed proof search without
turning a private operational log into an indiscriminate appendix.
\end{evidencebox}

This evidence does not make Git redundant. The trace shows what was tried; Git shows
what was accepted. Wotan explains the intended contract; project control explains
why work changed direction. A compiler error proves only that one intermediate term
failed. A later successful build and axiom audit establish the status of the accepted
endpoint.

\subsection{Contemporaneous knowledge and later interpretation}

Dates matter because the LaTeX source audit occurred after proof expansion closed.
Early work used the rendered PDF, text extraction, visual inspection, and local
mathlib experiments. The later source audit confirmed the closure convention,
polar sign, quantifiers, constants, and qualifications, but that confirmation was not
available to guide the first proofs. \ev{E01}\ev{E23}

Accordingly, phrases such as ``the source confirmed'' refer to retrospective
corroboration. Phrases such as ``the task record identified'' refer to evidence
available at the time. Interpretive claims -- for example, that interfaces acted as
epistemic firebreaks -- are explicitly presented as conclusions drawn from multiple
records, not as quotations from any one log.

\section{The initial problem and the first bet}
\label{sec:initial}

\subsection{A paper broader than a theorem list}

The repository began with a 60-page PDF and no Git or Wotan metadata. The paper mixed
linear-subspace geometry, closed convex cones, finite probability, temporal signal
analysis, nonnegative least squares, approximate embeddings, and a biological
interpretation. Its central mathematical contributions appeared mainly as numbered
equations and proof blocks rather than a conventional theorem hierarchy. The first
work therefore mapped equations, definitions, qualifications, and dependencies
rather than searching only for labels such as ``Theorem.'' \ev{E01}\ev{E26}

The feasibility review separated three questions:

\begin{enumerate}[label=\arabic*.]
  \item Do the mathematical conclusions follow from precise assumptions?
  \item Do idealized learning and population interfaces imply the claimed cone
    operations?
  \item Do actual neural populations satisfy those interfaces?
\end{enumerate}

Lean could address the first and conditional forms of the second. The third remained
empirical. This separation was not a final-report disclaimer added after success; it
was part of the first task contract and shaped the module boundary between general
mathematics and \artifact{Model/Population.lean}. \ev{E01}\ev{E05}

\subsection{The cautious forecast and the compressed calendar}

The initial report estimated four to eight person-weeks for a human Moreau and
double-rejection pilot and twelve to twenty-four person-months for broad Sections
2--5 coverage. Minutes later, the owner clarified that an autonomous AI agent was to
perform the work and expected hours rather than human project time. The response
reframed the four proposed work packages as Wotan idea tasks and selected the pilot
as the first decisive test. \ev{E01}\ev{E26}

The Git record from the feasibility commit through the source-audit commit spans
about five hours on 3 August 2026; the first scientific-report commit followed less
than half an hour later. This is a fact about preserved timestamps, not a controlled
productivity comparison. The original estimate assumed an experienced human,
maintainable mathlib-based work, and conventional review. The realized project used
an autonomous agent, aggressive scratch prototyping, a fixed local environment, and
a publication-oriented but repo-local theorem scope. It did not upstream a general
library or validate biology. \ev{E01}\ev{E25}\ev{E26}

\begin{turningpoint}[The first bet]
Instead of attempting the paper in order, the project chose one identity that
exercised all foundational decisions: metric projection, the source's non-positive
polar, Moreau decomposition, closure-safe cone operations, and the passage from a
pointwise equality to equality of closed conic hulls. If that identity failed under
the paper's quantification, the architecture would be reconsidered before broader
coverage. \ev{E01}\ev{E02}
\end{turningpoint}

\section{Contract before proof}
\label{sec:contract}

\subsection{Four ambiguities that could not remain informal}

The first permanent Lean code was preceded by
\artifact{FORMALIZATION-SPEC.md}. It fixed conventions whose omission would have
allowed an apparently successful proof to mean something different from the paper.
\ev{E02}

\begin{table}[htbp]
\centering
\caption{The initial theorem contract and its downstream consequences.}
\label{tab:initial-contract}
\begin{tabularx}{\linewidth}{@{}L{27mm}L{48mm}Y@{}}
\toprule
\textbf{Issue} & \textbf{Frozen interpretation} & \textbf{Why it mattered later} \\
\midrule
Closure &
Every arbitrary conic hull, sum, and image used by the cone algebra is closed in the
norm topology. &
Projected generator sets could not be treated as if their pointwise image were
already a closed cone; absorption and intersection required explicit hull arguments. \\
Polar sign &
The source polar is
\(C^\circ=\{z:\langle z,y\rangle\leq0\text{ for all }y\in C\}\), opposite in sign
to mathlib's ordinary dual-cone orientation. &
Moreau residuals, reflection laws, and rejection formulas could share one
paper-facing membership theorem without sign drift. \\
Quantification &
The cone core uses arbitrary real Hilbert spaces and arbitrary closed convex cones;
finite generation belongs to a separate population layer. &
The headline identities were not weakened to matrix or polyhedral special cases,
while finite active faces could later be introduced without contaminating the core. \\
Image equality &
A pointwise projection identity and equality of their closed conic hulls are separate
proof obligations. &
The double-rejection theorem first identifies individual projections and only then
lifts the result extensionally through \(\Conic\). \\
\bottomrule
\end{tabularx}
\end{table}

The Lean and mathlib versions were pinned at the same gate. Mathlib's Hilbert
projection theorem supplied existence of a minimizer, but no bundled cone-specific
Moreau API supplied the paper's exact bridge. A local probe demonstrated that a
noncomputable metric projection with membership and minimality could compile. The
session trace preserves even the correction of one mistaken singular-value API name
inside that probe. \ev{E01}\ev{E02}\ev{E26}

\subsection{Why coordinate-free geometry came first}

The specification deliberately avoided matrices in the core. Metric projection,
polarity, and orthogonality are the mathematical content of Moreau decomposition;
matrix coordinates are one representation. This choice later paid twice. The
subspace theory could be proved before a Moore--Penrose inverse existed, and the
finite population cones could live inside infinite-dimensional signal Hilbert
spaces because only their spans needed finite dimensionality. \ev{E02}\ev{E09}\ev{E15}\ev{E19}

This was not merely a preference for abstraction. It was a dependency decision. Had
the pilot been expressed through pseudoinverse matrices, the absence of a mathlib
Moore--Penrose API would have blocked the first gate for a reason unrelated to the
headline cone identity.

\section{The geometric pilot becomes an architecture}
\label{sec:pilot}

\subsection{Building the missing Moreau stack}

The pilot introduced a chosen metric projection, its variational characterization,
the polar residual, projection--residual orthogonality, Moreau decomposition,
bipolarity, and closure-safe conic projection and rejection. The final double-
rejection proof followed the paper's semantic architecture: establish a pointwise
orthogonal cone--polar decomposition, use Moreau uniqueness to identify the second
projection, and only then take closed conic hulls. \ev{E03}
This construction is the project-specific formal counterpart of the classical
polar-cone decomposition \cite{moreau1962,bauschke2017}.

The session trace shows a denser path than the commit. A failed rewrite in the metric
projection module, an unavailable local theorem identifier, a failed linear-
arithmetic closure, and an application mismatch in the operation layer occurred
within roughly three minutes. These were local elaboration and API repairs. They did
not force a change in the mathematical proof plan, which was already fixed by the
contract. \ev{E26}

A separate infrastructure failure did occur: dependency materialization was
interrupted by restricted network access. The recovery used the installed,
version-matched mathlib cache and obtained only the missing pinned dependency. The
resulting project retained a reproducible version manifest. The distinction is
important: network recovery changed how dependencies were materialized, not what the
theorem asserted. \ev{E03}\ev{E26}

\subsection{What passing the pilot established}

The pilot did more than prove one equation. It demonstrated that:

\begin{itemize}
  \item the source polar convention could be isolated cleanly;
  \item mathlib's Hilbert projection theorem was sufficient foundation for a local
    Moreau theory;
  \item arbitrary closed-cone quantification was viable;
  \item closure bookkeeping was manageable without hidden image-closedness
    assumptions; and
  \item the accepted endpoints could be built without \texttt{sorry} or custom
    axioms and audited at the declaration level.
\end{itemize}

These were the reasons the project continued, not merely the speed with which the
first commit appeared. Project control recorded the next gate as the much longer
intersection identity. \ev{E03}\ev{E04}

\subsection{Intersection and the first decomposition decision}

Before formalizing the intersection identity, the work considered whether polar
algebra might yield a shorter proof. The rejection-only form still required the
paper's four Moreau decompositions and cross-term control, so the direct geometry was
retained. One large theorem needed an increased heartbeat allowance, but the
kernel-checked term added no assumptions. \ev{E04}

At the same time, the approximate-invariance section was divided. The deterministic
projection-error implication was formalized with explicit relevant-vector and
transformed-optimality premises. The polarization step and sparse Gram bound were
independent and became a separate task. The random Johnson--Lindenstrauss guarantee
remained outside the deterministic theorem. This was the first repeated use of a
pattern that later dominated the project: prove the implication whose assumptions
are precise, name the missing bridge, and keep probability distinct from geometry.
\ev{E04}\ev{E06}

\section{Interfaces as epistemic firebreaks}
\label{sec:interfaces}

An explicit premise is sometimes mistaken for a failed formalization. In this
project, named interfaces had a more disciplined role. They prevented the proof
assistant from converting underspecified prose into hidden axioms while allowing
valid downstream consequences to be checked. Some interfaces remained model
boundaries; others were later discharged when a precise witness became available.

\begin{table}[htbp]
\centering
\caption{Representative interface life cycles.}
\label{tab:interface-life}
\begin{tabularx}{\linewidth}{@{}L{31mm}L{53mm}Y@{}}
\toprule
\textbf{Claim} & \textbf{Initial formal treatment} & \textbf{Later status} \\
\midrule
Population approximation (204) &
Named support-selected approximation premise in \artifact{PopulationRegion}. &
Remains a model premise; only (206)--(208) and zero-error hull consequences are
proved. No learning or biological dynamics is invented. \\
Stochastic NNLS convergence &
Recorded as an obligation because ``standard assumptions'' supplies no filtration,
sampling law, or moment bounds. &
Remains outside scope; deterministic batch convergence was developed instead. \\
Active-face formula (199) &
Initially a field of the population interface. &
Discharged by a finite sufficient KKT/Moreau certificate; necessity and a full
relative-interior partition remain unclaimed. \\
Compact convolution &
Boundary convention, finite-rank approximation, and compactness were separate named
premises. &
Discharged for a concrete nonzero causal unit-step Volterra operator, not for every
kernel or boundary convention. \\
Membrane \(O(\Delta t^2)\) term &
A named \artifact{QuadraticRemainder} interface. &
Discharged under an explicit Lipschitz drive with bound
\((L/\tau)|\Delta t|^2\). \\
Exponential occupancy approximation &
Represented by a named error-bound interface after the exact finite expectation. &
Discharged constructively for \(0<p\leq n\), with signed and absolute error bounds. \\
Two-message overlap &
Conditional on uniform marginals and targetwise independence. &
Exact for a later Cartesian-product experiment; still not a statement that
biological messages are independent. \\
\bottomrule
\end{tabularx}
\end{table}

The population layer made the boundary especially clear. Finite NNLS fixed points
and the error implications were provable. Support selection, learned approximation,
and biological adequacy were not. Recording the latter in
\artifact{MODEL-OBLIGATIONS.md} meant that later progress could target one obligation
without pretending the whole interface had been derived. \ev{E05}

\begin{figure}[htbp]
\centering
\begin{tikzpicture}[
  node distance=7mm and 13mm,
  every node/.style={font=\small,align=center},
  start/.style={draw=ReportGold,fill=ReportWarm,rounded corners,
    minimum width=37mm,minimum height=11mm},
  end/.style={draw=ReportBlue,fill=ReportPale,rounded corners,
    minimum width=38mm,minimum height=11mm},
  arrow/.style={-{Latex[length=2mm]},thick,draw=ReportGray}
]
  \node[start] (iface) {named interface};
  \node[end,right=of iface,yshift=15mm] (remain) {retained boundary\\equation (204), stochastic learning};
  \node[end,right=of iface] (special) {concrete witness\\causal Volterra operator};
  \node[end,right=of iface,yshift=-15mm] (derive) {general discharge\\membrane and occupancy bounds};
  \draw[arrow] (iface) -- (remain);
  \draw[arrow] (iface) -- (special);
  \draw[arrow] (iface) -- (derive);
\end{tikzpicture}
\caption{A named interface had three legitimate outcomes. Remaining explicit was not
an admission of proof; specialization and discharge required new mathematics.}
\label{fig:interfaces}
\end{figure}

The historical lesson is that an interface is useful only if it is visible,
source-shaped, and revisitable. A structure field whose name and contract match the
paper's approximation can be audited. A definition that simply says a population
``implements projection'' cannot.

\newpage
\section{Expansion by dependency, not by paper order}
\label{sec:expansion}

\subsection{The first empty queue}

The initial four tasks ended with a verified cone core, deterministic approximation
interfaces, finite NNLS fixed-point mathematics, and conditional population
consequences. The queue could have stopped there. Instead, a direction review compared
the achieved results with the paper-wide inventory and selected finite subspaces as
the next shared foundation. This was the first expansion decision made after the
owner's instruction to continue autonomously. \ev{E08}\ev{E09}

The chosen order did not follow the article. It followed dependency and risk:

\begin{enumerate}[label=\arabic*.]
  \item establish intrinsic finite-subspace operations and comparison measures;
  \item derive the NNLS step-size theorem from Gram bounds;
  \item close exact cone and circuit laws;
  \item formalize cone similarity and its finite-frame proxy;
  \item construct the exact finite occupancy experiment; and
  \item only then select the signal-space semantics needed by several analytic
    examples.
\end{enumerate}

This order allowed each task to leave a reusable interface for the next without
forcing an unresolved representation detail onto the critical path. \ev{E09}\ev{E10}\ev{E11}\ev{E12}\ev{E13}\ev{E14}

\begin{figure}[htbp]
\centering
\begin{tikzpicture}[
  node distance=7mm and 10mm,
  every node/.style={font=\scriptsize,align=center},
  core/.style={draw=ReportNavy,fill=ReportPale,rounded corners,
    minimum width=29mm,minimum height=10mm},
  branch/.style={draw=ReportCyan,fill=white,rounded corners,
    minimum width=29mm,minimum height=10mm},
  model/.style={draw=ReportGold,fill=ReportWarm,rounded corners,
    minimum width=29mm,minimum height=10mm},
  arrow/.style={-{Latex[length=1.8mm]},draw=ReportGray,thick}
]
  \node[core] (spec) {cone contract};
  \node[core,right=of spec] (moreau) {Moreau core};
  \node[core,right=of moreau] (exact) {exact cone laws};
  \node[branch,below=of spec] (subspace) {intrinsic subspaces};
  \node[branch,right=of subspace] (mp) {Moore--Penrose\\representation};
  \node[branch,below=of moreau] (nnls) {finite NNLS};
  \node[branch,right=of nnls] (rank) {spectral and\\rank-deficient convergence};
  \node[model,below=of subspace] (signal) {finite-window \(L^2\)};
  \node[model,right=of signal] (analytic) {wavelet, Volterra,\\membrane witnesses};
  \node[model,below=of nnls] (population) {active faces and\\population interface};
  \node[model,right=of population] (circuits) {conditional circuits};
  \draw[arrow] (spec) -- (moreau);
  \draw[arrow] (moreau) -- (exact);
  \draw[arrow] (subspace) -- (mp);
  \draw[arrow] (subspace) -- (nnls);
  \draw[arrow] (nnls) -- (rank);
  \draw[arrow] (moreau) -- (population);
  \draw[arrow] (population) -- (circuits);
  \draw[arrow] (signal) -- (analytic);
  \draw[arrow] (mp) -- (population);
  \draw[arrow] (rank) -- (population);
\end{tikzpicture}
\caption{A simplified dependency view. The historical sequence followed shared
foundations and exposed gaps rather than the linear order of paper sections.}
\label{fig:dependencies}
\end{figure}

\subsection{Intrinsic subspaces before pseudoinverse notation}

The subspace task proved the coordinate-free content of equations (12)--(23), the
range identity \(\Col(AA^\mathsf{T})=\Col(A)\), comparison formulas, and a concrete
failure of distributivity. Mathlib lacked the general Moore--Penrose API needed for
the paper's \(AA^+\) notation. Because none of those intrinsic results depended on
that notation, project control preserved the gap as a later task rather than blocking
the geometry. \ev{E09}

This split prevented two common distortions. The project did not assume a
pseudoinverse axiom merely to match notation, and it did not weaken the paper to
full-rank matrices. When the representation task was later executed, the construction
covered arbitrary rectangular, zero, and rank-deficient matrices. \ev{E15}

\subsection{The effective-map solution}

The final pseudoinverse did not use a full singular-value decomposition. It restricted
a linear map to the orthogonal complement of its kernel, corestricted it to its
range, proved the resulting effective map bijective, and inverted that equivalence.
The two products became the orthogonal projectors onto the range and kernel
complement; the four Penrose equations and uniqueness followed. \ev{E15}

The session record shows 25 successive patches to the permanent module. Several
failed rewrites were not failures of linear algebra but consequences of dependent
projection instances: equality of subspaces did not make the typeclass witnesses
definitionally identical. An extensional projector-equality lemma became the durable
bridge. \ev{E26}

This episode illustrates a recurring distinction. A coordinate-free theorem may be
mathematically simpler and more reusable, while a source-faithful coordinate formula
still has publication value. The appropriate architecture proves the former first
and derives the latter as a representation theorem.

\section{Failures at three different levels}
\label{sec:failures}

The polished Git history makes nearly every task appear to go from specification to
success in one commit. The session trace shows that this is a curated fiction useful
for version control but inadequate for history. Yet merely counting compiler errors
would be equally misleading. The failures fall into three levels.

\begin{table}[htbp]
\centering
\caption{Failure classes and their effects on the final development.}
\label{tab:failure-classes}
\begin{tabularx}{\linewidth}{@{}L{30mm}L{48mm}Y@{}}
\toprule
\textbf{Level} & \textbf{Typical evidence} & \textbf{Proper response} \\
\midrule
Elaboration and API &
Unknown theorem name, failed rewrite, coercion mismatch, stuck typeclass inference,
or a simplifier that sees the wrong normal form. &
Probe the library, expose intermediate expressions, replace a brittle tactic, or add
an extensional helper. The theorem contract normally remains unchanged. \\
Proof architecture &
Heartbeat timeout, unavailable abstraction, or a route that makes later obligations
intractable. &
Change the mathematical decomposition: use an effective map instead of SVD, direct
quadratic-form counting instead of an eigenvalue route, or simple-function density
instead of a custom time partition. \\
Theorem contract &
A counterexample, missing hypothesis, ambiguous boundary convention, or empirical
premise. &
Split the theorem, add the smallest visible premise, select a concrete model, or
leave the claim outside scope. Never disguise the repair as tactic work. \\
\bottomrule
\end{tabularx}
\end{table}

\subsection{Local failures: useful but not profound}

The sparse Gram prototype encountered unavailable summation lemmas, coercion
problems, and failed congruence applications over twelve patches. The NNLS prototype
encountered missing derivative combinators and rewrite failures over thirteen
patches. The wavelet scratch file accumulated 33 patches while interval integrals,
trigonometric normal forms, \(L^2\) representatives, and finite examples were made to
agree. These sequences matter because they show the actual cost of formal proof, but
none by itself refuted the source mathematics. \ev{E06}\ev{E08}\ev{E17}\ev{E26}

A better historical account records what those errors caused. In the Gram proof,
they encouraged an explicit expansion of diagonal and ordered off-diagonal terms,
which made the exact \(k-1\) count visible. In NNLS, they encouraged proving one
finite loss-increment identity from which both derivative and optimality results
could be derived. In the wavelet development, they produced reusable integer-mode
integral lemmas rather than a collection of numerical checks.

\subsection{Architectural failures: the route changes}

The active-face problem began with the paper's intrinsic language of faces and
relative interiors. Mathlib did not expose a matching polyhedral-face API. Rather
than build a full face lattice, the project formalized the exact sufficient
conditions used by the projection: a nonnegative active coefficient expansion,
polar inequalities against every generator, and orthogonality. Moreau uniqueness
then certified the selected span projection. Boundary certificates were allowed to
overlap, with equality of the projected point proved. This certificate-centered
view is consonant with the classical active-set treatment of nonnegative least
squares \cite{lawson1995}. \ev{E07}

This final theorem is both stronger and weaker than an informal equation (199). It is
stronger as a machine-checkable certificate: every finite inequality is visible and
degenerate active sets are covered. It is weaker as a classification: necessity and
a full relative-interior partition are not claimed. Calling this a ``workaround''
would miss the methodological point. It replaced an unavailable global theory with
the local theorem actually needed by the population interface.

The compact-convolution proof underwent an even clearer route change. A direct
finite-partition construction generated difficult elaboration and a deterministic
heartbeat timeout in \artifact{ScratchT18.lean}. The surviving proof introduced a
general map from an \(L^2\) Hilbert-valued kernel to a bounded analysis operator,
proved that indicator-constant kernels give rank-one operators, used density of
simple functions, and concluded compactness by norm closure. This was more general
than the planned special partition while remaining concrete about causal boundary
conditions. \ev{E18}\ev{E26}

\subsection{A mathematical obstruction: rank deficiency}

The most important failure was not a Lean error. For the projected batch NNLS map,
the familiar step condition \(0<\varepsilon<2/L\) does not imply a strict contraction
on coefficient space when the prediction map has a nontrivial kernel. A nonzero
kernel direction is fixed by the unprojected linear error map, so no factor \(q<1\)
can contract every coefficient difference. The project proved this obstruction
instead of hiding it behind a lower singular-value assumption. \ev{E10}

Two honest theorems replaced the hoped-for single statement:

\begin{itemize}
  \item under positive lower and upper spectral bounds, strict contraction gives a
    unique coefficient minimizer and geometric convergence;
  \item with only an upper bound and a nonempty minimizer set, Fejer descent,
    summable prediction error, finite-dimensional compactness, and a cluster
    minimizer give convergence to \emph{some} coefficient minimizer. All minimizers
    share the same prediction, while coefficients need not be unique.
\end{itemize}

The second proof was not a weakened version of the first. It was a different proof
architecture forced by the geometry of the kernel. Project control recorded this as
a redirection, and a later task completed the Fejer/compactness argument without a
positive lower spectral bound. \ev{E10}\ev{E16}

\begin{turningpoint}[Counterexamples are architectural evidence]
The kernel obstruction did more than close a false route. It determined the public
API: prediction uniqueness and coefficient uniqueness became separate theorems, and
the rank-deficient convergence module could not accidentally acquire a full-rank
hypothesis. \ev{E10}\ev{E16}
\end{turningpoint}

\section{The analytic arc: semantics before calculation}
\label{sec:analysis-arc}

\subsection{Choosing what a signal is}

The signal layer could not begin from a convolution formula alone. It first selected
Lebesgue measure restricted to \([0,T]\) and Bochner \(L^2\) classes for scalar and
finite-vector signals. This decision immediately corrected two informal habits.
Pointwise matrix identities had to be stated almost everywhere because \(L^2\)
elements are equivalence classes. And the image of a closed cone under a bounded
linear map could not simply be declared closed; the formal cone image used a closed
conic hull. \ev{E14}

The first task at this layer stopped deliberately at named interfaces for boundary
conditions, generic convolution realization, compactness, Laplace transfer, and the
quadratic membrane remainder. Exact linearity, norm bounds, conic aggregation, and
finite-dictionary facts were proved. This left a coherent theorem layer without
pretending that unspecified analytic hypotheses had been supplied by notation.
\ev{E14}

\subsection{The wavelet normalization false alarm}

During the initial PDF review, extracted text made the normalization factor for the
three Hann-windowed wavelets look suspicious. The session update records a genuine
concern that the displayed constant might not give norm one. High-resolution
inspection of the rendered formula resolved the scope of the square root, and later
exact integration proved the source constant correct. \ev{E01}\ev{E17}\ev{E26}

This episode is useful because it shows two distinct error channels. PDF text
extraction can corrupt notation; a numerical or symbolic plausibility check can then
raise the right question without settling it. The final answer required an exact
Lean proof in the chosen recording-window measure.

The proof history was substantial. \artifact{ScratchT17.lean} first probed missing
integral theorem names, then derived integer sine and cosine integrals, product-to-sum
identities, Hann-square expansions, means, Gram entries, and \(L^2\) inner products.
A deterministic elaboration timeout and many local mismatches were resolved before
the file was renamed and moved into
\artifact{IPNPCNS/Examples/OrthonormalWavelets.lean}. Its final theorems also checked
three NNLS examples and kept baseline subtraction separate from literal nonnegative
firing rates. \ev{E17}\ev{E26}

\subsection{A concrete compact filter, not a generic slogan}

The source's compact-filter discussion did not fix every kernel extension or boundary
convention. The formalization chose a nonzero unit-step causal response on the finite
window, giving the Volterra operator
\[
  (Vx)(t)=\int_0^t x(s)\,ds.
\]
The final module proves its almost-everywhere integral formula, the bound
\(\lVert V\rVert\leq T\), finite-rank approximation in operator norm, and
compactness. It separately proves that the chosen realization is not periodic and is
not the alternative whole-window zero-extension formula. \ev{E18}

The missing Hilbert--Schmidt abstraction therefore did not become an axiom. It became
a prompt to prove compactness from simpler certified pieces. The resulting theorem
is narrower in kernel scope but stronger in semantic clarity than an unqualified
claim about ``the convolution operator.''

\subsection{The membrane remainder and the value of a modest constant}

Equation (33)'s \(O(\Delta t^2)\) term originally remained behind a named
\artifact{QuadraticRemainder}. An empty-queue review later selected it as a bounded
analytic gap. The exact first-order low-pass step was written as an exponentially
weighted Bochner integral. Under a Lipschitz premise on the net drive, subtracting
the frozen initial drive produced an explicit remainder bounded by
\((L/\tau)|\Delta t|^2\). \ev{E20}

A sharper elementary factor of one half was considered. The final proof used the
uniform interval bound instead, because it supplied the required step-independent
quadratic constant with a smaller and more auditable argument. This is a small but
representative scientific choice: formal verification need not optimize every
constant when the claimed order and interface are preserved, but it must say which
constant was actually proved. \ev{E20}\ev{E26}

\section{The probabilistic arc: construct the experiment}
\label{sec:probability-arc}

The occupancy formula after equation (30) appears elementary, but its meaning depends
on a sampling model. The formal experiment made each of \(m\) axons choose uniformly
from the \(p\)-element subsets of \(n\) targets. Choices are without replacement
within one axon and independent between axons. Treating all \(mp\) collaterals as
independent target draws would have proved a different formula and was explicitly
rejected. \ev{E13}

The exact expected occupied-target count followed by indicator expansion and finite
product averaging. Zero targets, zero axons, zero collaterals, and the impossible
case \(p>n\) received separate statements. The source's exponential replacement and
two-message overlap were initially left behind different interfaces because they
used different extra assumptions. \ev{E13}

Later empty-queue reviews returned to both interfaces. For \(0<n\) and \(p\leq n\),
elementary exponential remainder and finite power-difference bounds proved that the
exponential approximation is a lower bound with absolute count error at most
\(mp^2/n\). This was a genuine discharge of an error interface. \ev{E21}

The overlap statement required a different move. Two messages were sampled from the
literal Cartesian product of the exact one-message sample space. Targetwise
independence then followed from the construction, and expected overlap became exactly
\(q^2/n\), where \(q\) is the exact one-message expectation. Independence was made
structural in the selected mathematical experiment; it was not inferred for neural
messages. \ev{E22}

The session trace reveals the final proof repair. Directly unfolding the structure
predicates failed. Changing the goals to the literal finite-expectation equalities
made the product factorization visible, after which the existing abstract overlap
theorem could be instantiated. \ev{E26}

\section{Steering as part of the proof method}
\label{sec:steering}

Wotan and project control are subjects of this report because they influenced which
mathematics was attempted. They are not presented as scientific evidence for the
theorems themselves. Lean's kernel and the audited source tree play that role.

\subsection{What Wotan contributed}

Each task combined a claim-oriented context, acceptance criteria, hard verification
endpoints, manual checks, risks, an approach, implementation notes, obstacles,
evidence, and an outcome. The initial continuations were stored as inactive ideas.
After the owner authorized continuous autonomous work, project control promoted them
into an executable queue and permitted creation of follow-on tasks. \ev{E01}\ev{E02}

This mattered most at boundaries. The sparse Gram proof was split from the already
verified projection-error theorem. Active-face construction and deterministic NNLS
convergence were separated from the conditional population interface. The missing
pseudoinverse notation was preserved without delaying intrinsic subspaces or NNLS.
The source audit and two reports were blocked behind a moving proof-phase barrier so
that they would describe a stable theorem set. \ev{E04}\ev{E05}\ev{E09}\ev{E23}

The logs also make negative scope durable. A theorem not attempted because its
stochastic assumptions were absent does not silently disappear; it remains in
\artifact{MODEL-OBLIGATIONS.md} and in the task outcome. This is especially important
for interdisciplinary work, where an omitted empirical bridge can otherwise be
mistaken for a proved consequence.

\subsection{What project control contributed}

Project-control reviews occurred at scientific gates and whenever the actionable
queue became empty. Each review recorded the current evidence, uncertainty, selected
direction, and revisit condition. The pattern was not ``always continue.'' It included
preserve, evaluate, redirect, and eventually close proof expansion in favor of source
conformance. \ev{E03}\ev{E09}\ev{E18}\ev{E22}\ev{E23}

Three types of steering decision were especially productive:

\begin{description}
  \item[Protect the critical path.]
  Prove coordinate-free geometry before pseudoinverse representation; separate the
  already valid deterministic implication from an independent probability theorem.
  \item[Change proof architecture when the mathematics demands it.]
  Preserve the full-rank contraction theorem and create a distinct rank-deficient
  Fejer task after the kernel obstruction.
  \item[Reopen named boundaries selectively.]
  When the queue emptied, choose a compact, source-shaped obligation such as span
  reduction, the membrane remainder, or occupancy error rather than expanding into
  an underspecified stochastic or biological program.
\end{description}

The final redirection is as important as the expansions. After the independent
two-message theorem, the remaining open areas required a new stochastic or
biological specification, generic transform theory, or a much larger intrinsic face
development. Project control judged source fidelity more limiting than another small
proof and released the LaTeX audit. \ev{E22}\ev{E23}

\begin{turningpoint}[Autonomy requires visible stopping logic]
Continuous autonomous execution did not mean that every adjacent mathematical topic
became authorized scope. The combination of a maintained obligation register, an
empty-queue review, and explicit revisit conditions supplied a reasoned boundary
between productive continuation and unbounded theorem accumulation.
\end{turningpoint}

\section{The verification standard grew with the project}
\label{sec:verification-standard}

``It compiles'' meant different things at different moments. The first feasibility
probe established only that mathlib could support a chosen metric projection. The
final closure claim combined a complete build, proof-gap and unsafe-declaration
scans, public theorem-signature review, a 38-endpoint axiom audit, equation-level
source conformance, and report artifact QA. Treating these checks as interchangeable
would erase how the standard of evidence matured. \ev{E01}\ev{E03}\ev{E23}\ev{E24}

\begin{table}[htbp]
\centering
\caption{Layers of evidence added during the project.}
\label{tab:verification-layers}
\begin{tabularx}{\linewidth}{@{}C{11mm}L{39mm}Y@{}}
\toprule
\textbf{Layer} & \textbf{Check} & \textbf{Question answered} \\
\midrule
1 & Focused probe or scratch compilation &
Does the proposed type, library bridge, or proof route elaborate in the pinned
environment? \\
2 & Focused module build &
Does the accepted module compile with its actual imports and namespace? \\
3 & Complete \artifact{lake build} &
Does the whole dependency graph elaborate together? The final build completed 2,960
jobs. \\
4 & Gap and declaration scan &
Do project sources contain \texttt{sorry}, a custom \texttt{axiom}, or an
\texttt{unsafe} declaration? The final scans found none. \\
5 & Signature and \texttt{\#print axioms} audit &
Does the public theorem quantify as intended, and on which logical principles does
the compiled constant depend? \\
6 & LaTeX source conformance &
Were signs, constants, domains, quantifiers, equation references, and nearby
qualifications read correctly from the authoritative source? \\
7 & PDF extraction and visual QA &
Does the reader-facing artifact contain the intended text, metadata, tables,
figures, and pagination without visual defects? \\
\bottomrule
\end{tabularx}
\end{table}

Lean's small trusted kernel checks elaborated proof terms; mathlib provides previously
checked definitions and theorems \cite{demoura2015lean,mathlib2020}. The project adds
no custom axiom. The representative endpoints report only
\texttt{propext}, \texttt{Classical.choice}, and \texttt{Quot.sound}, the accepted
standard principles used by the chosen Lean/mathlib development. This audit does not
mean that the code is constructive or computational: metric projection and several
finite-dimensional choices are intentionally noncomputable. \ev{E23}

\subsection{Verification itself had correction loops}

The audit process was not exempt from failure. A task-level axiom check once ran in
parallel before the new object files existed. The source-wide audit similarly raced
the full build and initially reported a missing wavelet object file. Five temporary
signal endpoint names were also misspelled in the first audit script. Sequential
reruns with corrected names passed; no project theorem changed. \ev{E05}\ev{E23}\ev{E26}

These incidents justify preserving commands and outputs rather than recording only
``audit passed.'' They also show why verification evidence should be reproducible and
why a failed audit command must be classified before it is interpreted as a failed
theorem.

The build-job count grew from 2,356 at the foundational specification to 2,960 at
proof closure. The number is an environment-specific Lake build count, not a direct
measure of project lines or theorem difficulty. Its historical value is that every
closed task rechecked an expanding dependency graph rather than compiling only the
new file. \ev{E02}\ev{E23}

\section{The source audit as an independent closure gate}
\label{sec:source-audit}

The proof work initially relied on the distributed PDF, extracted text, rendered
pages, and manual cross-checks. The owner later required a stronger gate: after proof
expansion was judged complete, inspect the author-supplied LaTeX archive, reopen any
affected proof, and only then write the two reports. Project control deliberately
placed this gate after the last compact theorem task, so it would audit a stable
claimed scope rather than a moving target. \ev{E22}\ev{E23}

The archive audit recorded SHA-256 provenance, inventoried one TeX root and twelve
figures, expanded source macros, compiled the article through three passes, and
compared the resulting 60-page text with the distributed PDF. The body matched; the
observed difference was the arXiv title-page stamp and induced layout. The
source-to-Lean trace then classified claims through equation (210) as exact,
conditional, constructively strengthened, represented, or outside scope. \ev{E23}

Two findings required correction. A wavelet module header referred to Section 5
instead of Section 3.2. The transmission-invariance boundaries around equations
(158)--(173) were not collected in the obligation document. Both were documentation
defects. No formula, theorem signature, or proof had to be reopened. \ev{E23}

\begin{keyresult}[Meaning of the clean source audit]
The audit supports the claim that the finished Lean endpoints have the stated
relationship to the source. It does not claim line-by-line formalization of every
citation, empirical assertion, or research direction in the paper
\cite{nilsson2026}. A clean conformance result validates the declared boundary; it
does not erase that boundary.
\end{keyresult}

The ordering matters historically. The source audit retrospectively confirmed that
the early closure, polar, and Hilbert-space decisions were faithful. It cannot be
used as evidence that no uncertainty existed while those decisions were first made.

\section{What the session record reveals between commits}
\label{sec:between-commits}

The commit series contains one accepted state per task. The session trace reveals a
different rhythm: propose a small term, compile, inspect the exact goal, patch, and
compile again. High-risk analytic tasks were often developed in scratch files and
moved only after the public names and proofs stabilized. \ev{E25}\ev{E26}

\begin{longtable}{@{}L{37mm}C{18mm}L{91mm}@{}}
\caption{Selected scratch histories preserved outside Git.}
\label{tab:scratch-history}\\
\toprule
\textbf{Scratch path} & \textbf{Patches} & \textbf{What the trace adds to the final account} \\
\midrule
\endfirsthead
\toprule
\textbf{Scratch path} & \textbf{Patches} & \textbf{What the trace adds to the final account} \\
\midrule
\endhead
\path{/private/tmp/gram-proto.lean} & 12 &
Shows the progression from unavailable sum lemmas and coercion failures to the
explicit ordered-pair count that preserved \(k-1\). \\
\path{/private/tmp/nnls-proto.lean} & 13 &
Shows derivative-API attempts, the exact loss-increment pivot, KKT decomposition,
and the later convergence layer before permanent modules were created. \\
\artifact{ScratchT16.lean} & 6 &
Shows the local Fejer/compactness construction used after strict contraction became
mathematically impossible. \\
\artifact{ScratchT17.lean} & 33 &
Shows the entire exact-integration and \(L^2\) route, including unavailable library
names, timeouts, Gram calculations, and final move to the wavelet module. \\
\artifact{ScratchT18.lean} & 26 &
Shows the abandoned elaboration-heavy route and the change to a general
kernel-analysis operator plus simple-function density. \\
\artifact{ScratchT20.lean} & 15 &
Shows exact exponential-kernel integration, algebraic normalization failures, and
the final Lipschitz remainder theorem. \\
\artifact{ScratchT25.lean} & 9 &
Shows why direct structure unfolding was replaced by literal finite-expectation
goals for the product law. \\
\bottomrule
\end{longtable}

The patch count is not a difficulty score. A long diff may be a mechanical rename;
a one-line counterexample may force a new theorem architecture. The counts are useful
only to refute the impression that the clean modules appeared in one pass and to
locate where the dense proof search occurred.

The trace also retains actions that Git intentionally omits: temporary axiom-audit
files, theorem-name probes, failed builds, report renderings, and scratch cleanup. In
that sense it is closer to a laboratory notebook than the commit graph. Unlike a
laboratory notebook, it is machine-generated and exhaustive only for tool-mediated
events. The human-readable session audit selects evidence without publishing raw
operational content. \ev{E26}

\section{Scale, speed, and what they do not prove}
\label{sec:scale}

At proof closure the Lean tree contained 31 module files under
\artifact{IPNPCNS/}, 8,367 Lean lines including the root import module, 439 theorem or
lemma declarations, and 172 definitions, structures, classes, or abbreviations. The
Git series had 23 commits through source audit and 24 through the scientific report.
The permanent mathematical modules ranged from 20 lines for a finite example wrapper
to more than 1,000 lines for the exact wavelet development. \ev{E25}

The elapsed commit window is striking. The first feasibility commit is timestamped
17:31:59 local time, the source-audit commit 22:36:58, and the first report commit
23:05:38, all on the same date. The live session itself began earlier at 15:10:21Z
and continued through report production. \ev{E25}\ev{E26}

Three cautions prevent this from becoming a misleading performance claim.

\begin{enumerate}[label=\arabic*.]
  \item Commit timestamps bound accepted repository states; they do not measure
    uninterrupted active compute, equivalent human labor, or independent review.
  \item The initial estimate concerned experienced human development and included a
    broad uncertainty margin. It was not a benchmark prediction for an autonomous
    agent operating in a pinned, cached environment.
  \item The final scope is broad but disciplined. It does not formalize every cited
    background theorem, generic stochastic learning, all convolution kernels, a full
    intrinsic polyhedral-face theory, or empirical CNS adequacy.
\end{enumerate}

The honest conclusion is therefore not that formal mathematics has become
instantaneous. It is that an autonomous agent can compress a large amount of
repo-local formal proof work when the environment, task contracts, and feedback
loops are favorable. The compression makes evidence discipline more important, not
less: without the session trace, clean commits would conceal the dense search; without
source and axiom audits, speed could be mistaken for assurance.

\section{General lessons for interdisciplinary formalization}
\label{sec:lessons}

The project suggests several lessons that extend beyond this paper.

\subsection{Formalize semantic operations before convenient representations}

The coordinate-free cone and subspace layers survived every later extension.
Pseudoinverse matrices, finite generators, and population circuits could be added as
representation theorems. Starting from source notation alone would have coupled the
core to missing library infrastructure. \ev{E02}\ev{E09}\ev{E15}

\subsection{Treat quantifiers and closure as part of the result}

Changing from arbitrary closed cones to finite polyhedral cones would have made some
proofs easier but weakened the headline result. Omitting norm closure would have made
several image arguments shorter but false as stated. These are not implementation
details; they determine which theorem was verified. \ev{E02}\ev{E03}\ev{E11}

\subsection{Separate exact mathematics, approximation, and empirical realization}

Near-isometry implications, random embedding probabilities, learned population
accuracy, and biological validity occupy different logical levels. The formalization
benefited from different namespaces and interfaces for those levels. A conditional
theorem can be strong and useful without converting its premise into an empirical
fact. \ev{E04}\ev{E05}\ev{E23}

\subsection{Let counterexamples shape the public API}

The rank-deficient kernel obstruction forced prediction and coefficient uniqueness
apart. Empty cones forced semantic theorems to carry nonemptiness hypotheses. The
nonattainment of an infinite-dimensional infimum led to a lower-bound proof rather
than selection of a minimizer. These are improvements in mathematical statement,
not setbacks in implementation. \ev{E10}\ev{E12}\ev{E16}

\subsection{Use interfaces as revisitable hypotheses}

An interface should name the missing mathematical content narrowly enough that it
can later be discharged or specialized. The membrane, exponential occupancy, and
convolution interfaces passed this test. Equation (204) also passes it by remaining
visible: the project can prove consequences without claiming a training theorem it
does not possess. \ev{E05}\ev{E18}\ev{E20}\ev{E21}

\subsection{Concrete witnesses can be more honest than generic abstractions}

The Volterra theorem did not establish every square-integrable convolution claim.
It established one nontrivial causal compact operator with explicit boundary
semantics. The product occupancy theorem did not establish biological independence;
it established the exact result for a chosen independent experiment. Such witnesses
turn vague existence prose into auditable mathematics without overgeneralizing.
\ev{E18}\ev{E22}

\subsection{Proof logs and proof artifacts answer different questions}

Git and compiled constants answer what survived. Session patches and diagnostics
answer how the proof was found. Wotan and project control answer what the work was
supposed to establish and why it changed direction. A credible history needs all
three, with provenance labels that prevent one from impersonating another. \ev{E26}

\subsection{Audit the source after the theorem set stabilizes, but before publication}

An early source audit would have chased a moving coverage boundary. A post-publication
audit would have been too late to correct theorem or report claims. The selected gate
made every material source discrepancy capable of reopening proof work while keeping
the traceability exercise finite. \ev{E22}\ev{E23}

\subsection{Verification is layered assurance}

Compilation, absence of admitted proofs, axiom inspection, theorem-signature review,
source conformance, and visual report QA answer different questions. None subsumes
the others. The project became more credible as these layers accumulated and their
failures were themselves recorded. \ev{E23}\ev{E24}\ev{E26}

\subsection{Autonomous speed increases the value of explicit control}

When proof production is fast, scope can expand faster than reflective judgment. An
obligation register, decision gates, and a rule for empty queues kept the work tied
to source-shaped value. The mechanism did not make the decisions correct by fiat; it
made their evidence and uncertainty inspectable.

\section{Residual uncertainty and limits of this history}
\label{sec:limits}

The report is evidence-based, not omniscient.

First, Wotan and project-control records are contemporaneous but self-authored. They
summarize intent and interpretation; they are not independent observations. Git and
compiler output corroborate implementation claims, while the source audit
corroborates final conformance.

Second, the session trace is unusually complete for tool-mediated editing. It does
not record unexpressed cognition, and this report does not expose model-internal
reasoning. The checksum fixes the prefix actually audited; later session events are
outside that historical snapshot. \ev{E26}

Third, the 38-endpoint axiom audit is a broad representative sample across all
mathematical layers, not a printed list for all 439 theorem and lemma declarations.
The repository-wide scan separately establishes the absence of admitted proofs,
custom axioms, and unsafe declarations. The full build checks dependency closure.
\ev{E23}

Fourth, the source audit validates the project's claims about its theorem set. It
does not formalize every cited result or resolve the deliberately open stochastic,
generic analytic, intrinsic-face, network-stability, or empirical questions. The
proof corpus should be described as a substantial verified mathematical layer, not
as a complete formal theory of the central nervous system. \ev{E23}

Finally, no independent external Lean reviewer or upstream mathlib maintainer is part
of the recorded project. Kernel checking gives strong logical assurance; independent
review could still improve names, abstractions, proof maintainability, and the
scientific choice of interfaces.

\section{Conclusion}
\label{sec:conclusion}

The IPNPCNS formalization did not emerge by mechanically transcribing a paper. It
emerged by exposing the commitments hidden in mathematical prose and allowing each
one to meet a different kind of pressure: type theory, library availability,
counterexample, analytic regularity, sampling semantics, biological uncertainty, or
source conformance.

The decisive moves were architectural. The Moreau pilot fixed a coordinate-free
closed-cone core. Named interfaces kept underspecified learning and approximation
claims honest. A kernel counterexample split full-rank and rank-deficient NNLS.
Sufficient active certificates replaced an unavailable global face theory. An
effective-map inverse supplied Moore--Penrose representation without a rank
assumption. Exact Fourier work, finite-rank kernel approximation, Lipschitz analysis,
and finite product experiments converted selected analytic and probabilistic
interfaces into theorems.

The historical record is as important as the sequence of successes. Git preserves
the accepted checkpoints; Wotan and project control preserve scope and decisions;
the full session trace preserves scratch code and compiler-guided search that the
final tree intentionally omits. Their combination makes it possible to say not only
that the development now builds, but how its verification boundary was discovered.

The most general lesson is simple: a formalization is credible when its theorem
contracts, failed routes, assumptions, trusted base, and source relationship are all
visible. Autonomous proof production can be very fast. The evidentiary work that
makes the result scientifically legible cannot be skipped.

\clearpage
\appendix

\section{Chronological spine of accepted states}
\label{app:chronology}

The main report is organized by causal arcs. This appendix supplies the minimal
chronological spine needed to date those arcs. Times are local commit times on
3 August 2026. A row identifies an accepted Git state, not all work between states.

\begin{longtable}{@{}C{13mm}L{17mm}L{39mm}L{72mm}@{}}
\caption{Git checkpoints through the scientific report.}
\label{tab:git-spine}\\
\toprule
\textbf{Time} & \textbf{Commit} & \textbf{Checkpoint} & \textbf{Historical role} \\
\midrule
\endfirsthead
\toprule
\textbf{Time} & \textbf{Commit} & \textbf{Checkpoint} & \textbf{Historical role} \\
\midrule
\endhead
17:31 & \commit{e291e52} & Feasibility &
Paper inventory, Lean probe, empirical boundary, and four inactive continuation
ideas. \\
17:46 & \commit{eda50fa} & Specification &
Pinned environment, closure and polar conventions, and exact pilot contract. \\
17:53 & \commit{87b1d51} & Moreau pilot &
Metric projection, Moreau decomposition, cone operations, and double rejection. \\
18:10 & \commit{7a29065} & Intersection &
Three rejection formulas and deterministic projection-error theorem. \\
18:17 & \commit{f0db876} & Population boundary &
Finite NNLS fixed points, explicit approximation interface, and conditional
consequences. \\
18:30 & \commit{084cf8f} & Approximation constants &
Polarization and direct sparse Gram bound. \\
18:38 & \commit{d5f9d34} & Active certificate &
Finite sufficient active-face region and population constructor. \\
19:00 & \commit{529e111} & NNLS optimality &
Loss derivative, KKT/minimizer equivalence, and contraction-based convergence. \\
19:18 & \commit{560c81e} & Subspaces &
Intrinsic finite geometry, matrix range identity, and comparison measures. \\
19:26 & \commit{48efc13} & Spectral NNLS &
Full-rank contraction factor and rank-deficient obstruction. \\
19:36 & \commit{10baf10} & Cone laws and circuits &
Exact identities, reflection, primitive circuits, and conditional control. \\
19:46 & \commit{d6f610a} & Cone comparison &
Worst-case similarity, finite-frame proxy, and coverage error. \\
19:56 & \commit{a7ecddb} & Occupancy &
Exact one-message finite experiment and visible approximation boundaries. \\
20:09 & \commit{b26dcec} & Signal semantics &
Finite-window Bochner \(L^2\), bounded filters, dictionaries, and analytic
interfaces. \\
20:32 & \commit{764a00f} & Moore--Penrose &
Effective-map inverse, Penrose equations, uniqueness, and projector renderings. \\
20:42 & \commit{94c4bfa} & Rank-deficient NNLS &
Fejer/compactness convergence to a selected minimizer. \\
21:15 & \commit{af5d096} & Exact wavelets &
Hann-window normalization, orthonormality, examples, and baseline convention. \\
21:46 & \commit{08b72d3} & Compact Volterra filter &
Causal realization, norm bound, finite-rank approximation, and compactness. \\
21:54 & \commit{a0a608e} & Projection structure &
Positive homogeneity and finite-span reduction. \\
22:05 & \commit{3dd6ff9} & Membrane remainder &
Exact low-pass step and Lipschitz quadratic error. \\
22:15 & \commit{123d09a} & Occupancy approximation &
Signed exponential comparison and explicit \(mp^2/n\) error. \\
22:21 & \commit{367a1c8} & Product overlap &
Independent two-message experiment and exact \(q^2/n\) expectation. \\
22:36 & \commit{fd1361b} & Source audit &
LaTeX conformance, traceability through equation (210), and final trust checks. \\
23:05 & \commit{5a83b13} & Scientific report &
Separate 24-page non-chronological mathematical account and PDF QA. \\
\bottomrule
\end{longtable}

\section{Evidence ledger}
\label{app:evidence-ledger}

The entries below are intentionally compact. Full task evidence appears in
\artifact{wotan/dev-log/}; steering context appears in
\artifact{PROJECT-CONTROL.md}; source mapping appears in
\artifact{SOURCE-AUDIT.md}; and the session extraction method appears in
\artifact{reports/session-trace-audit.md}.

\begin{description}[leftmargin=16mm,labelwidth=12mm,style=multiline]
\evidenceitem{E01}{Contemporaneous feasibility report and task log,
\path{LEAN-FEASIBILITY.md} and \path{wotan/dev-log/T-0001.md}, preserved by
commit \commit{e291e52}. Supports the initial state, scope distinction, risk
inventory, pilot choice, and human-effort forecast.}

\evidenceitem{E02}{\artifact{FORMALIZATION-SPEC.md}, project-control records 001,
and commit \commit{eda50fa}. Supports the pinned environment, closure and polar
conventions, arbitrary-Hilbert-space target, and explicit non-goals.}

\evidenceitem{E03}{Task log T-0003 and commit \commit{87b1d51}. Supports the Moreau
API, pointwise-before-hull proof structure, network-cache recovery, and first
assumption/axiom audit.}

\evidenceitem{E04}{Task log T-0004, project-control reviews 002--003, and commit
\commit{7a29065}. Supports the direct four-Moreau intersection architecture,
heartbeat adjustment, deterministic near-isometry scope, and decomposition of
independent endpoints.}

\evidenceitem{E05}{Task log T-0005, \artifact{MODEL-OBLIGATIONS.md}, project-control
review 004, and commit \commit{f0db876}. Supports separation of NNLS mathematics,
population assumptions, stochastic claims, and empirical validity.}

\evidenceitem{E06}{Task log T-0006 and commit \commit{084cf8f}. Supports the sharp
polarization constant and the direct active quadratic-form proof of the sparse Gram
bound.}

\evidenceitem{E07}{Task log T-0007, project-control reviews 005--006, and commit
\commit{d5f9d34}. Supports the sufficient KKT/Moreau active certificate, treatment of
boundaries and the empty face, and the unclaimed intrinsic-face converse.}

\evidenceitem{E08}{Task log T-0008, project-control reviews 006--007, and commit
\commit{529e111}. Supports exact loss expansion, KKT/global-minimum equivalence,
contraction-based convergence, and the early recognition of the rank-deficient
boundary.}

\evidenceitem{E09}{Task log T-0009, project-control reviews 007--008, and commit
\commit{560c81e}. Supports dependency-led expansion, intrinsic subspace results, and
deferral of the pseudoinverse representation layer.}

\evidenceitem{E10}{Task log T-0010, project-control review 009, and commit
\commit{48efc13}. Supports the full-rank contraction factor, common prediction, and
the proved kernel obstruction to strict coefficient contraction.}

\evidenceitem{E11}{Task log T-0011 and commit \commit{10baf10}. Supports exact cone
laws, reflection and circuit rewrites, explicit control premises, and rejection of an
invalid image-membership shortcut.}

\evidenceitem{E12}{Task log T-0012 and commit \commit{d6f610a}. Supports explicit
nonemptiness, nonattainment-safe infimum reasoning, optimistic finite proxies, and
coverage-conditional error.}

\evidenceitem{E13}{Task log T-0013, project-control reviews 011--012, and commit
\commit{a7ecddb}. Supports the exact without-replacement/independent-between-axons
experiment, boundary cases, and separate approximation interfaces.}

\evidenceitem{E14}{Task log T-0014, project-control reviews 012--013, and commit
\commit{b26dcec}. Supports finite-window Bochner \(L^2\), almost-everywhere semantics,
closed image hulls, and the initial analytic interface boundary.}

\evidenceitem{E15}{Task log T-0015, project-control reviews 013--014, and commit
\commit{764a00f}. Supports the effective-map Moore--Penrose construction, all-rank
scope, four Penrose equations, and dependent projector-equality issue.}

\evidenceitem{E16}{Task log T-0016, project-control reviews 014--015, and commit
\commit{94c4bfa}. Supports the rank-deficient Fejer, compactness, cluster-minimizer,
and full-sequence convergence proof without a lower spectral bound.}

\evidenceitem{E17}{Task log T-0017, project-control reviews 015--016, and commit
\commit{af5d096}. Supports exact wavelet constants, local Fourier integration,
orthonormality, example coefficients and residuals, and baseline convention.}

\evidenceitem{E18}{Task log T-0018, project-control reviews 016--017, and commit
\commit{08b72d3}. Supports the simple-function/rank-one compactness route, causal
boundary choice, norm bound, and separation of alternative boundaries.}

\evidenceitem{E19}{Task log T-0019, project-control reviews 017--018, and commit
\commit{a0a608e}. Supports empty-queue selection, positive homogeneity, span
reduction, unrestricted ambient Hilbert space, and dependent-instance rewrite repair.}

\evidenceitem{E20}{Task log T-0020, project-control reviews 018--019, and commit
\commit{3dd6ff9}. Supports the exact low-pass integral, Lipschitz premise, selected
quadratic constant, and decision not to pursue the sharper half-factor.}

\evidenceitem{E21}{Task log T-0021, project-control reviews 019--021, and commit
\commit{123d09a}. Supports constructive discharge of the exponential occupancy error
interface and separate invalid regimes.}

\evidenceitem{E22}{Task log T-0025, project-control reviews 021--022, and commit
\commit{367a1c8}. Supports the Cartesian-product two-message experiment, exact
overlap expectation, and proof-phase closure decision.}

\evidenceitem{E23}{\artifact{SOURCE-AUDIT.md}, task log T-0022, project-control
reviews 022--023, and commit \commit{fd1361b}. Supports archive provenance, source
build and comparison, equation-level traceability, documentation corrections, final
2,960-job build, gap scan, and 38-endpoint axiom audit.}

\evidenceitem{E24}{Task log T-0023, project-control review 024, and commit
\commit{5a83b13}. Supports the separate scientific-report role, 24-page artifact,
and render-correct QA loop.}

\evidenceitem{E25}{Git log and current Lean-tree measurement. Supports the accepted
checkpoint order and timestamps, 31 module files, 8,367 Lean lines, 439 theorem or
lemma declarations, and 172 definition-like declarations. Timestamps do not measure
human-equivalent effort.}

\enlargethispage{2\baselineskip}
\evidenceitem{E26}{\RaggedRight \artifact{reports/session-trace-audit.md} and the main Codex
session prefix with SHA-256
{\footnotesize\texttt{80e1b5977d4ff07fbfc696f319eb2a31}\allowbreak
\texttt{db33d57636464e20b81d1063ba763ba5}}.
Supports the scratch-file reconstruction, 427 successful patch events, 174 changes
over 33 temporary or audit files, compiler-feedback chains, and report QA history.}
\end{description}

\clearpage
\begin{thebibliography}{9}

\bibitem{nilsson2026}
M. N. P. Nilsson.
\newblock Information processing by neuron populations in the central nervous
system: A theory of the mathematical structure of data and operations.
\newblock arXiv:2309.02332v3, 2026.

\bibitem{moreau1962}
J.-J. Moreau.
\newblock Decomposition orthogonale d'un espace hilbertien selon deux cones
mutuellement polaires.
\newblock \emph{C. R. Acad. Sci. Paris}, 255:238--240, 1962.

\bibitem{lawson1995}
C. L. Lawson and R. J. Hanson.
\newblock \emph{Solving Least Squares Problems}.
\newblock SIAM, Philadelphia, 1995.

\bibitem{bauschke2017}
H. H. Bauschke and P. L. Combettes.
\newblock \emph{Convex Analysis and Monotone Operator Theory in Hilbert Spaces}.
\newblock Springer, second edition, 2017.

\bibitem{demoura2015lean}
L. de Moura, S. Kong, J. Avigad, F. van Doorn, and J. von Raumer.
\newblock The Lean theorem prover (system description).
\newblock In \emph{Automated Deduction -- CADE-25}, pages 378--388.
Springer, 2015.

\bibitem{mathlib2020}
The mathlib Community.
\newblock The Lean mathematical library.
\newblock In \emph{Proceedings of the 9th ACM SIGPLAN International Conference on
Certified Programs and Proofs}, pages 367--381, 2020.

\bibitem{lakatos1976}
I. Lakatos.
\newblock \emph{Proofs and Refutations: The Logic of Mathematical Discovery}.
\newblock Cambridge University Press, 1976.

\end{thebibliography}

\end{document}
